\section{Limitations}\label{sec:equiv}

A potential limitation of our technique is that the \transDname{} produces
nondeterministic code that does not reduce \emph{only} to the result
of demand semantics in general. Some logic languages have a built-in facility to collapse a
nondeterministic computation into a deterministic one, returning either the
first result found or a list of all possible results. In languages that also
give the user control over the order in which branches are tried, we can arrange
for less-defined results to always be tried first, thereby ensuring that we
always obtain the minimal result. We take this approach in
\cref{sec:case-study}.

As discussed in \cref{sec:case-study}, our technique is somewhat slow with even
modestly-sized inputs. Although unfortunate, we consider this to be
acceptable. A major motivation behind this line of research was the desire to
find a lightweight representation of demand functions that could be used to test
their asymptotic bounds (constants included) \emph{before} trying to formally
verify them. For this purpose, it usually suffices to test quite small inputs,
where our technique performs quite well in absolute (wall-clock) terms.
